\documentclass[11pt]{article}

\usepackage[margin=1.1in]{geometry}
\usepackage{amsmath,amssymb,amsthm,mathtools}
\usepackage{booktabs}
\usepackage{longtable}
\usepackage{enumitem}
\usepackage[colorlinks=true,linkcolor=blue,citecolor=blue]{hyperref}
\usepackage{xcolor}

\newcommand{\PROVEN}{\textcolor{teal}{\textsc{[proven]}}}
\newcommand{\DEFINITION}{\textcolor{gray}{\textsc{[definition]}}}
\newcommand{\CONJECTURAL}{\textcolor{orange}{\textsc{[conjectural]}}}
\newcommand{\BLOCKED}{\textcolor{red}{\textsc{[blocked]}}}

\newcommand{\PT}{\Pi}
\newcommand{\PS}{\mathcal{P}(\PT)}
\newcommand{\WS}{\mathcal{W}}
\newcommand{\fib}[1]{F_{#1}}
\newcommand{\kerEq}{\equiv_{K}}
\newcommand{\Lin}{\mathrm{Lin}}
\newcommand{\Hser}{H_{\mathrm{ser}}}
\newcommand{\dW}{d_{\mathcal{W}}}
\newcommand{\dP}{d_{\mathcal{P}}}
\newcommand{\R}{\mathbb{R}}
\newcommand{\N}{\mathbb{N}}
\newcommand{\Z}{\mathbb{Z}}

\theoremstyle{plain}
\newtheorem{theorem}{Theorem}[section]
\newtheorem{conjecture}[theorem]{Conjecture}
\newtheorem{lemma}[theorem]{Lemma}
\theoremstyle{definition}
\newtheorem{definition}[theorem]{Definition}
\newtheorem{example}[theorem]{Example}
\theoremstyle{remark}
\newtheorem{remark}[theorem]{Remark}
\newtheorem{ladder}[theorem]{Ladder Note}

\title{The Mathematics of Multi Fractal Workflow:\\
One Object, Every Level\\[0.5em]
\large A complete mathematical ascent --- from counting and statistics to
optimal transport, stratified geometry, multifractal spectra, and
machine-checked foundations --- organized around a single map
$\tau : \mathcal{P}(\Pi) \to \mathcal{W}$}
\author{Sean Chatman}
\date{July 14, 2026}

\begin{document}
\maketitle

\begin{abstract}
\noindent Multi Fractal Workflow (MFW) studies one map: the transformation
$\tau$ that compiles PDDL~3.1 planning behavior into POWL~v2 workflow
structure. What makes MFW unusual is that this single object is
simultaneously a statistics problem (a partition of a sample space), a
combinatorics problem (counting linear extensions of partial orders), an
algebra problem (a quotient by an equivalence relation), a logic problem (a
machine-checked theory in dependent type theory), an information-theory
problem (conditional entropy of a lossy code), a measure-theory and
optimal-transport problem (geometry on conditional fiber measures), a
differential-geometry problem (dimension loss on a stratified space), a
functional-analysis problem (a vector space of invariant observables), a
fractal-geometry problem (a bundle of R\'enyi dimensions), and an
algebraic-topology problem (a simplicial complex of concurrency). This
paper walks the entire ladder in order, rung by rung, defining every piece
of mathematics as it is first needed and showing exactly where in MFW it
lives. Every statement carries its verification standing from the Lean~4
formalization (\texttt{ProcInt.MFW}): \textsc{proven},
\textsc{definition}, \textsc{conjectural}, or \textsc{blocked}. Nothing is
assumed beyond one college statistics course; nothing relevant is left out
at the top.
\end{abstract}

\tableofcontents

% ================================================================
\section{The one object}
\label{sec:object}

Everything in this paper is about a single map and its preimages:
\[
  \tau : (\PS,\ \dP,\ \nu) \longrightarrow (\WS,\ \dW,\ \tau_{*}\nu).
\]
Here $\PS$ is the set of \emph{lawful behaviors} of a PDDL~3.1 planning
theory $\PT$ (timed action sequences that obey the physics, achieve the
goal, and satisfy every temporal, numeric, and trajectory constraint), and
$\WS$ is the set of POWL~v2 \emph{workflow objects} (hierarchical models
built from partial orders, which carry concurrency, and choice graphs,
which carry decisions and cycles). The map $\tau$ is the compiler: behavior
in, workflow class out. The decorations $\dP, \nu, \dW, \tau_{*}\nu$ ---
a metric, a measure, and their images --- are exactly the higher mathematics
this paper climbs toward.

The central open problem, stated once now and revisited at every rung:

\begin{conjecture}[Crown theorem]\CONJECTURAL\label{conj:crown}
$\tau(b_1) = \tau(b_2) \iff b_1 \sim_I b_2$, where $\sim_I$ is
Mazurkiewicz trace equivalence under the action-independence relation $I$.
Equivalently: $\ker(\tau) = \mathrm{TraceEq}_I$.
\end{conjecture}

Each level of mathematics below is a different pair of glasses for looking
at this one statement and its consequences.

\begin{ladder}
The paper is a ladder. Each section states (a) the mathematics it
introduces, at the level of a good first course in that subject; (b) the
MFW objects built from it; (c) what is proven versus conjectured. Later
sections freely use earlier ones.
\end{ladder}

% ================================================================
\section{Rung 0: Sets, functions, counting (the statistics you have)}
\label{sec:rung0}

\subsection{Mathematics}
A \emph{set} is a collection; a \emph{function} $f : X \to Y$ assigns each
element of $X$ one element of $Y$. $f$ is \emph{injective} if distinct
inputs give distinct outputs, \emph{surjective} if everything in $Y$ is
hit, \emph{bijective} if both. A \emph{partition} of $X$ is a family of
nonempty, pairwise-disjoint subsets covering $X$ --- exactly the ``grouping
variable'' of statistics. The number of orderings of $k$ distinct items is
$k!$; the binomial coefficient $\binom{n}{2}$ counts pairs. A finite
probability distribution on $X$ assigns $p(x) \ge 0$ with $\sum_x p(x) = 1$.

\subsection{In MFW}
\begin{definition}[Fiber]\DEFINITION\;
$\fib{w} = \tau^{-1}(w) = \{b \in \PS : \tau(b) = w\}$ --- the group of
behaviors the compiler merges into workflow $w$.
\end{definition}

\begin{theorem}[Fibers partition the sample space]\PROVEN\;
Every $b$ lies in $\fib{\tau(b)}$ (\texttt{fiber\_partition}), and distinct
fibers are disjoint: $w_1 \ne w_2 \implies \fib{w_1} \cap \fib{w_2} =
\emptyset$ (\texttt{fiber\_disjoint}).
\end{theorem}

Already at Rung 0 a design was falsified: a proposed workflow similarity
used the Jaccard index $J(\fib{w_1}, \fib{w_2}) =
|\fib{w_1} \cap \fib{w_2}| / |\fib{w_1} \cup \fib{w_2}|$. By disjointness
$J \equiv 0$ off the diagonal, so the induced ``metric'' is the useless
discrete one. A proven set-theory triviality killed a geometry proposal
before implementation; the repair arrives at Rung 8.

Counting enters through fibers too:
\begin{definition}[Fiber cardinality and probability]\DEFINITION\;
For a finite enumeration of behaviors,
$|\fib{w}|$ is a count, and $p(w) = |\fib{w}| / |\PS|$ is the fraction of
behaviors landing in $w$ --- a genuine probability distribution on the image
of $\tau$ ($p \ge 0$: \PROVEN).
\end{definition}

% ================================================================
\section{Rung 1: Relations, equivalence, and quotients (discrete algebra)}
\label{sec:rung1}

\subsection{Mathematics}
A \emph{binary relation} on $X$ is a set of pairs. An \emph{equivalence
relation} $\equiv$ is reflexive ($x \equiv x$), symmetric, and transitive;
its \emph{equivalence classes} $[x] = \{y : y \equiv x\}$ form a partition,
and conversely every partition arises this way. The \emph{quotient}
$X/{\equiv}$ is the set of classes. The \emph{kernel} of a function $f$ is
the equivalence $x \equiv_f y \iff f(x) = f(y)$; the first isomorphism
idea: $f$ factors as $X \twoheadrightarrow X/{\equiv_f} \hookrightarrow Y$.
A \emph{strict partial order} $\prec$ is irreflexive and transitive; a
\emph{linear extension} of $\prec$ on $n$ items is a total order
compatible with it, formally a bijection $\sigma$ with
$i \prec j \implies \sigma(i) < \sigma(j)$.

\subsection{In MFW}
The kernel of the compiler is the star of the show:
\begin{definition}[KernelEquiv]\DEFINITION\;
$b_1 \kerEq b_2 \iff \tau(b_1) = \tau(b_2)$. Equivalence laws: \PROVEN.
\end{definition}

Independence and traces give the \emph{other} equivalence:
\begin{definition}[Independence; Mazurkiewicz trace equivalence]\DEFINITION\;
$I \subseteq A \times A$ is symmetric; $(a,b) \in I$ demands five
commutation facts (effects commute; preconditions stable; invariants
stable; numeric flows compatible; trajectory constraints preserved).
Trace equivalence $u \sim_I v$ on action sequences is the
reflexive--symmetric--transitive closure of the adjacent swap
$u\,a\,b\,v \leadsto u\,b\,a\,v$ for $(a,b) \in I$. That $\sim_I$ is an
equivalence relation: \PROVEN{} (\texttt{traceEquiv\_equivalence}).
\end{definition}

The crown conjecture (\ref{conj:crown}) now reads as pure Rung-1 algebra:
\emph{two independently defined equivalence relations on the same set are
claimed equal}. The formalization proves the claim is non-circular
(\PROVEN): $\kerEq$ unfolds using only $\tau$ and equality in $\WS$;
$\sim_I$ unfolds using only $I$ and event lists --- neither side mentions
the other's vocabulary (\texttt{kernelEquiv\_unfolds\_without\_traces},
\texttt{traceEquiv\_unfolds\_without\_tau}, both \texttt{rfl}).

Partial orders enter as the \emph{content} of a trace class:
\begin{definition}[Causal order; trace class]\DEFINITION\;
Each behavior $b$ under $I$ induces a causal order $P_b$ ($i \prec j$ iff
dependent and earlier; irreflexive, transitive). The trace class is
$[b]_I = \{b' : b \sim_I b'\}$.
\end{definition}

\begin{theorem}[Fiber = trace class, conditionally]\PROVEN\ under hypothesis\;
If the crown holds (taken as an explicit hypothesis), then
$\fib{\tau(b)} = [b]_I$ as sets (\texttt{fiber\_eq\_traceClass}).
\end{theorem}

\begin{conjecture}[Trace class $\simeq$ linear extensions]\CONJECTURAL\;
Under bijection hypotheses (distinct events, decidable $I$, independence
exactly complementing transition-level dependence):
$[b]_I \simeq \Lin(P_b)$. Classical (Mazurkiewicz 1977; Diekert--Rozenberg
1995) but not yet formalized under PDDL~3.1 semantics.
\end{conjecture}

There is also an \emph{order on the transformations themselves}:
$\tau_1 \sqsubseteq \tau_2$ (``refines'') iff $\tau_1$-equality implies
$\tau_2$-equality, i.e.\ $\tau_2 = \psi \circ \tau_1$; this is reflexive
and transitive (\PROVEN) and will matter at Rung 6.

% ================================================================
\section{Rung 2: Graphs, Petri nets, and structural decomposition}
\label{sec:rung2}

\subsection{Mathematics}
A \emph{directed graph} is nodes plus arrows. A \emph{Petri net}
$N = (P, T, F)$ is a bipartite directed graph of \emph{places} and
\emph{transitions}; tokens on places move by \emph{firing} transitions
(consume one token per input place, produce one per output place).
A \emph{workflow net} has a unique source place, unique sink place, and
every node on a path between them; it is \emph{safe} if no place ever
holds two tokens and \emph{sound} if there are no dead transitions, every
reachable state can still finish, and completion is clean. Two structural
subclasses: \emph{state machines} (every transition: $\le 1$ input,
$\le 1$ output place --- all branching is choice) and \emph{marked graphs}
(every place: $\le 1$ input, $\le 1$ output transition --- all branching is
concurrency).

\subsection{In MFW and its companion conversion theory}
POWL~v2 is defined recursively (Kourani--van Zelst):
\begin{definition}[Powl model, well-formedness]\DEFINITION\;
\[
  \mathrm{Powl}(\alpha) ::= \mathrm{atom}(a) \mid \mathrm{silent} \mid
  \mathrm{xor}(\vec{c}) \mid \mathrm{loop}(c_1, c_2) \mid
  \mathrm{po}(\vec{c}, \prec),
\]
with hereditary well-formedness ($|\vec{c}| \ge 2$ for xor; $\prec$
irreflexive and transitive for po). POWL v2 adds \emph{choice graphs}
(directed start-to-end graphs over submodels, expressing non-block
decisions and cycles); a POWL v2 object carries the model, its
well-formedness proof, hierarchical depth, boundary variables
$\partial w$, and choice graphs.
\end{definition}

The structural duality --- partial order $\leftrightarrow$ marked graph,
choice graph $\leftrightarrow$ state machine --- powers the proven
conversion theory for the companion question (Kourani--Park--van der
Aalst 2026): \emph{separable} WF-nets (recursively nested state machines
and marked graphs, never cross-linked) convert to POWL~2.0
\emph{completely} (language equality, their Thm.~5.5) and \emph{totally on
the class} (their Thm.~5.6); explicit sound free-choice counterexamples
(long-term dependencies; loops entering one parallel branch; interleaved
choice/concurrency at a shared synchronization) bound the class from
outside. In MFW's terms this settles the \emph{existence} and
\emph{exactness} legs of the compilation question; the \emph{maximality}
leg is the crown conjecture.

% ================================================================
\section{Rung 3: Combinatorics of order --- counting concurrency}
\label{sec:rung3}

\subsection{Mathematics}
For a finite poset $P$, $e(P) = |\Lin(P)|$ counts linear extensions ---
between $1$ (a chain) and $n!$ (an antichain). An \emph{antichain} is a set
of pairwise-incomparable elements; the \emph{width} $w(P)$ is the largest
antichain, and Dilworth's theorem says $w(P)$ equals the minimum number of
chains covering $P$. \emph{Interval orders}: if each event occupies a real
interval $[s_a, e_a]$, then $a \prec b \iff e_a < s_b$.

\subsection{In MFW}
\begin{definition}[Serialization entropy]\DEFINITION\;
$\Hser(P) = \log e(P)$: the information content of accidental
serialization. One concurrency class of $k$ mutually independent actions
carries $\Hser = \log k!$ --- the compression opportunity of the
transformation is \emph{factorial}.
\end{definition}

\begin{definition}[Erasure target]\DEFINITION\;
The least causal relation preserving correctness:
$R^{*} = \arg\min_R |R|$ subject to $\Lin(E, R) \subseteq
\mathrm{Valid}(\PT)$. ``Maximize concurrency'' \emph{means} emit $R^{*}$.
\end{definition}

\begin{definition}[Width; temporal abstraction]\DEFINITION\;
$w(P)$ is the primary concurrency coordinate (one scalar --- deliberately
too weak; its multiscale refinement is Rung 9). The abstraction ladder
$\mathrm{MetricTime} \to \mathrm{IntervalTime} \to \mathrm{CausalTime} \to
\mathrm{HierarchicalTime}$ erases timestamps in stages; interval orders
sit in the middle. The governing slogan, recorded in the source:
\emph{the question determines the sufficient abstraction}.
\end{definition}

\begin{definition}[Temporal restriction gap]\DEFINITION\;
With a simple temporal network $N$ (difference bounds
$\ell_{ij} \le t_j - t_i \le u_{ij}$), $\Lin_T(P,N)$ is the temporally
feasible extensions, $H_T(P,N) = \log|\Lin_T(P,N)|$
(\texttt{opaque} interface in Lean), and
$\Hser(P) - H_T(P,N) \ge 0$ measures causal concurrency the clock forbids.
Four concurrency types are kept distinct:
$C_E = C_C \cap C_T \cap C_R$ (executable = causal $\cap$ temporal
$\cap$ resource).
\end{definition}

% ================================================================
\section{Rung 4: Logic and type theory --- the machine-checked substrate}
\label{sec:rung4}

\subsection{Mathematics}
In \emph{dependent type theory} (the logic of Lean~4), propositions are
types and proofs are their inhabitants. A \emph{subtype}
$\{x : X \mathrel{/\!/} P(x)\}$ packages an element with a proof, so
``being lawful'' can be made impossible to fake rather than merely
asserted. Definitional equality (\texttt{rfl}) distinguishes unfoldings
from theorems. \emph{Decidability} of a relation is itself a mathematical
asset (it licenses computation inside proofs).

\subsection{In MFW}
This rung explains several audit repairs that are really logic:
\begin{itemize}[nosep]
  \item \textbf{Proof-carrying behaviors.} The old
  \texttt{LawfulBehavior} stored lawfulness as data fields, so a behavior
  with \texttt{achievesGoal := False} inhabited the ``lawful'' space. The
  repair is the subtype discipline: \texttt{BehaviorTrace} (raw data)
  $+$ \texttt{IsLawful} (predicate, five clauses) $+$
  \texttt{LawfulBehavior} (the pair). \emph{The proof is a witness, not a
  payload.}
  \item \textbf{Semantics or absence.} A constraint type without a
  satisfaction relation does not constrain anything; hence $\PT$ carries
  $\models_T, \models_N, \models_\Gamma$ explicitly.
  \item \textbf{No vacuous laws.} A former field
  \texttt{respects\_equivalence} concluded $\mathrm{True}$ from an
  assumption already containing its conclusion; it was deleted, and the
  kernel characterization became a \emph{theorem target} instead of an
  assumption --- the difference between logic and wishful typing.
  \item \textbf{Non-circularity as \texttt{rfl}.} Both crown sides unfold
  definitionally in disjoint vocabularies; the biconditional is therefore
  contentful.
  \item \textbf{Standing discipline.} Every declaration carries
  \textsc{proven} / \textsc{definition} / \textsc{conjectural} /
  \textsc{blocked}; this paper inherits that labeling.
\end{itemize}

% ================================================================
\section{Rung 5: Information theory --- what the compiler erases}
\label{sec:rung5}

\subsection{Mathematics}
Shannon entropy of a uniform distribution on $m$ items is $\log m$;
conditional entropy $H(X \mid Y) = \sum_y p(y) H(X \mid Y = y)$ measures
what remains unknown about $X$ after seeing $Y$; conditional mutual
information $I(X; Y \mid Z) \ge 0$ measures shared information given $Z$.
A statistic $T$ is \emph{sufficient} for $Y$ when $I(Y; X \mid T) = 0$
(Fisher--Neyman).

\subsection{In MFW}
The compiler is a lossy code; the loss is exactly quantified:
\begin{definition}[Fiber entropy; total fiber entropy]\DEFINITION\;
\[
  S_\tau(w) = \log|\fib{w}|, \qquad
  H(B \mid W) = \sum_w p(w)\, S_\tau(w), \qquad
  0 \le H(B \mid W) \le \log|\PS|.
\]
Nonnegativity, the zero-iff-singleton characterization
($S_\tau(w) = 0 \iff |\fib{w}| \le 1$; for reachable $w$, injectivity),
monotonicity in fiber size, the total-collapse identity
($S = \log|\PS|$ when everything maps to one class), and the upper bound
all have complete Lean proof terms but carry in-source \CONJECTURAL{}
labels pending API review; congruence and $|\fib{w}| > 0$ for reachable
$w$ are \PROVEN.
\end{definition}

\begin{remark}[Additivity repair]
$\log|\fib{w}|$ is not additive over disjoint fibers; the additive atom is
$\eta(w) = p(w) \log|\fib{w}|$, and $H(B \mid W) = \sum_w \eta(w)$
distributes correctly over hierarchies. This became the \emph{seventh}
measure kind after an audit found fiber entropy conflated with temporal
slack --- they are orthogonal: $\mathrm{Slack}(w) \ne H(B \mid W = w)$.
\end{remark}

\begin{conjecture}[Entropy collapse]\CONJECTURAL\;
If the crown and the trace-class bijection both close:
\[
  S_\tau(\tau(b)) = \log|\fib{\tau(b)}| = \log|[b]_I|
  = \log|\Lin(P_b)| = \log e(P_b) = \Hser(P_b).
\]
Fiber entropy \emph{equals} serialization entropy: an information-theoretic
quantity (Rung 5) collapses onto a combinatorial one (Rung 3), and the
erased information becomes computable from the emitted diagram alone.
\end{conjecture}

Sufficiency completes the rung: $\tau$ is \emph{sufficient for} a target
property $Y$ when $Y$ factors through $\tau$ (deterministic fragment,
\DEFINITION); the stochastic bridge $I(Y; B \mid W) = 0 \implies$
sufficiency is \CONJECTURAL{} (placeholder structure
\texttt{SemanticResidualInfo}, awaiting measure-theoretic conditional
entropy --- Rung 8's machinery).

% ================================================================
\section{Rung 6: Set-level functional structure --- observability and
sigma-algebras}
\label{sec:rung6}

\subsection{Mathematics}
A \emph{sigma-algebra} on $X$ is a family of subsets closed under
complement and countable union; the sigma-algebra \emph{generated by} a map
$f$ is $\{f^{-1}(U)\}$ --- the events decidable by observing $f$ only.
A \emph{Boolean algebra} of propositions is closed under
$\neg, \wedge, \vee, \to$. A \emph{vector space} of functions is closed
under pointwise $+$ and scalar $\cdot$; a subset \emph{spans} if finite
linear combinations reach everything; a family is a \emph{basis} if it
spans and is linearly independent.

\subsection{In MFW}
\begin{definition}[Observable / hidden / fiber-constant]\DEFINITION\;
$P : \PS \to \mathrm{Prop}$ is \emph{observable} iff it factors:
$\exists \hat P,\ P = \hat P \circ \tau$; \emph{fiber-constant} iff
$\tau(b_1) = \tau(b_2) \implies (P(b_1) \leftrightarrow P(b_2))$;
\emph{hidden} iff some fiber contains a $P$-witness and a
$\neg P$-witness.
\end{definition}

\begin{theorem}[The characterization and its algebra]
Observable $\iff$ fiber-constant (forward \PROVEN; reverse constructed
via the factoring predicate, labeled \CONJECTURAL{} in-source though the
term closes). Hidden $\iff$ not fiber-constant; observable
$\implies$ not hidden (\PROVEN). Observables form a Boolean algebra
containing $\top, \bot$ and all pullbacks, closed under every unary and
binary connective (\PROVEN).
\end{theorem}

\begin{definition}[Observable sigma-algebra; information horizon]\DEFINITION\;
$\mathcal{O}_\tau = \{S \subseteq \PS : S \text{ saturated}\}
= \{\tau^{-1}(U) : U \subseteq \WS\}$ --- both inclusions \PROVEN. The
\emph{information horizon} $H_\tau$ is the observable/hidden dichotomy;
inside it, everything is determinable from the POWL object; outside,
erased. Refinement $\tau_1 \sqsubseteq \tau_2$ monotonically enlarges the
observable algebra (\PROVEN) --- the principled dial for ``make $Y$
observable, and no more.''
\end{definition}

\begin{definition}[Invariant observables $F_\tau$]\DEFINITION\;
$F_\tau = \{f : \PS \to \R \mid f \text{ constant on fibers}\}$. It
contains $0$ and constants; closed under $+$, scalar, negation, and
product (\PROVEN); every $g \circ \tau$ belongs (\PROVEN); conversely
every $f \in F_\tau$ factors through surjective $\tau$ by choosing fiber
representatives (term closes; \CONJECTURAL{} label in-source). Hence
morally $F_\tau \cong (\WS \to \R)$; for finite $\WS$,
$\dim F_\tau = |\mathrm{im}\,\tau|$. This is the collider-physics analogy
made exact: $F_\tau$ plays the role of infrared/collinear-\emph{safe}
observables, blind to what the transformation erases.
\end{definition}

\begin{definition}[Basis program; truncation]\DEFINITION\;
Five primitive kinds (causal incidence, interval, motif, hierarchical,
temporal-boundary --- EFP analogues: angular structure, energy flow, graph
polynomials, jet substructure, thrust). Span $\subseteq F_\tau$: complete
induction proof (labeled \CONJECTURAL). The basis claim ---
$\mathrm{span}\{\varphi_i\} = F^{\mathrm{adm}}_\tau$ under an admission
profile --- is \CONJECTURAL, with conjectured truncation bounds
$\#\mathrm{causal} \le \binom{\mathrm{maxConc}}{2}$,
$\#\mathrm{motif} \le 2^{\binom{\mathrm{maxBranch}}{2}}$, total order
polynomial in the profile. Feature extraction
$b \mapsto [\varphi_1(b), \dots, \varphi_n(b)]$ is fiber-constant
(\PROVEN): equivalent behaviors get identical feature vectors.
\end{definition}

% ================================================================
\section{Rung 7: Probability refined --- the q-lens and R\'enyi reweighting}
\label{sec:rung7}

\subsection{Mathematics}
Given a positive finite distribution $p$, the \emph{escort} (or
R\'enyi/thermodynamic) deformation at inverse-temperature-like parameter
$q \in \R$ is $p^{(q)}(x) = p(x)^q / \sum_y p(y)^q$. As $q$ grows it
concentrates on the mode; $q = 1$ returns $p$; $q = 0$ is uniform on the
support; $q < 0$ inverts attention onto rare outcomes.

\subsection{In MFW}
\begin{theorem}[q-lens]\PROVEN\;
$L_q(p)$ is a well-defined distribution (positivity, normalization), and
the ratio law holds:
\[
  \frac{L_q(p)(x)}{L_q(p)(y)} = \left(\frac{p(x)}{p(y)}\right)^{q}
  \qquad (\texttt{qLens\_ratio\_law}).
\]
\end{theorem}
Interpretation fixed by the source: $q$ is \emph{attention control over
behavioral mass inside POWL choice structure} --- $q \gg 0$ audits dominant
channels, $q \ll 0$ audits the rare-behavior tail, $q = 1$ is the Shannon
scale. The q-lens also parameterizes \emph{search rails} (Rung 10's
portfolio) and, at Rung 9, the multifractal spectrum.

% ================================================================
\section{Rung 8: Measure theory and optimal transport --- geometry on
workflows}
\label{sec:rung8}

\subsection{Mathematics}
A \emph{measure} $\nu$ assigns sizes to sets; the \emph{pushforward}
$\tau_*\nu$ transports it along a map ($\tau_*\nu(U) =
\nu(\tau^{-1}U)$); a \emph{conditional measure} $\nu_w =
\nu(\cdot \mid \tau = w)$ restricts and renormalizes to a fiber. A
\emph{metric} satisfies identity, symmetry, triangle inequality
(\emph{pseudo}-metric: distinct points may sit at distance 0). Given a
ground metric $d$, the \emph{Wasserstein distance} between distributions is
the optimal-transport cost
$W_p(\mu, \rho) = \inf_{\pi \in \Pi(\mu,\rho)} \mathbb{E}_\pi[d(x,y)]$,
infimum over couplings with the given marginals (Villani).

\subsection{In MFW}
This is the repair of the Rung-0 Jaccard collapse:
\begin{definition}[Wasserstein workflow metric]\DEFINITION\;
Fix a behavior metric $\dP$ on $\PS$ (axioms carried as proof fields).
Define
\[
  \dW(w_1, w_2) = W_p(\nu_{w_1}, \nu_{w_2}).
\]
Fibers are disjoint as \emph{sets}, but their behavior \emph{distributions}
can be close in $\dP$ --- distance lives at the level of measures, not
supports. In Lean the distance is \texttt{opaque} (interface; the transport
construction needs measure theory beyond the module). A
\texttt{WorkflowMetric} record induces a \texttt{PseudoMetricSpace}
(\CONJECTURAL{} pending the distance's own well-definedness).
\end{definition}

\begin{definition}[Pushforward masses; provenance law]\DEFINITION\;
\texttt{PushforwardMass}: $\mu : \WS \to \R_{\ge 0}$. The full
\emph{vector measure} has seven components
$[\mu_B, \mu_T, \mu_C, \mu_L, \mu_{Sl}, \mu_F, \mu_S]$ (behavioral,
temporal, choice, linearization, slack, fluent, entropic), with the
standing \emph{provenance requirement}: each must ultimately be derived
from $(\tau, \nu)$, not posited.
\end{definition}

Rung 8 also supplies what Rung 5 lacked: measure-theoretic conditional
entropy $H(B \mid W = w)$ generalizing $\log|\fib{w}|$ to infinite fibers,
and the machinery for the sufficiency bridge $I(Y; B \mid W) = 0$ --- both
explicitly deferred as research targets in the source.

% ================================================================
\section{Rung 9: Fractal geometry --- the multifractal spectrum bundle}
\label{sec:rung9}

\subsection{Mathematics}
For a measure spread over a nested family of partitions at scales $k$, the
\emph{partition sum} $Z(q,k) = \sum_{B \in P_k} p_B^q$ generates the
\emph{R\'enyi (generalized) dimensions} $D_q$ from its scaling in $k$;
$D_0$ is box-counting (capacity), $D_1$ the information dimension, $D_2$
correlation. A measure is \emph{multifractal} when $D_q$ genuinely varies
with $q$ --- different moments see different scaling. $Z$ is classically
log-convex in $q$.

\subsection{In MFW}
POWL v2 supplies the canonical multiscale filtration for free: the
hierarchy $P_0 \preceq P_1 \preceq \cdots \preceq P_h$ of the workflow
object (refinement law currently concludes \texttt{True} --- the
is-submodel-of predicate is \CONJECTURAL).

\begin{definition}[Spectrum bundle]\DEFINITION\;
\[
  \vec{D}_q = \big(D_q^{B}, D_q^{T}, D_q^{C}, D_q^{L}, D_q^{Sl},
  D_q^{F}\big),
\]
one generalized dimension per measure kind: a section of a trivial
$\R^6$-bundle over the $q$-line, with the axiom that $D_0$ is finite for
each kind. \emph{Provenance (audit):} the bundle must be
\emph{manufactured} from admitted measure + scale + fit evidence, never
constructed as an arbitrary function --- without the kernel, ``spectrum''
is a name, not a derived object.
\end{definition}

\begin{definition}[Cross-spectrum tension]\;
$\Delta_{ij}(q) = D_q^i - D_q^j$. Antisymmetry, diagonal vanishing, zero
trace: \PROVEN. The interpretation dictionary: $D_q^B \gg D_q^T$ temporal
bottleneck; $D_q^B \ll D_q^T$ scheduling elasticity; $D_q^C \gg D_q^B$
choice sparsity; $D_q^L \gg D_q^T$ \emph{false concurrency} --- Rung 3's
temporal restriction gap reappearing as a spectral signature. The source is
explicit: \emph{the optimization signal is the tension profile, not any
single spectrum}.
\end{definition}

\begin{definition}[Partition sums; scale resolution]\;
$Z(q,k) = \sum_{B \in P_k} p_B^q$ with $Z(0,k) = |P_k|$ (term closes,
requires positivity; \CONJECTURAL{} label), $Z(1,k) = 1$ under
normalization (\PROVEN), nonnegativity and disjoint additivity (terms
close; \CONJECTURAL{} labels). The scale-resolved bundle $\vec{D}_q(k)$
adds hierarchical depth; tension antisymmetry persists at every depth
(\PROVEN). Open: log-convexity of $Z$ in $q$; whether the bundle is a
genuine fiber bundle with interesting structure group; a Fisher-metric
reading of the profile space.
\end{definition}

\begin{definition}[Transformation information profile]\DEFINITION\;
The master diagnostic
$I_\tau = (d_{\mathrm{int}}, \vec{D}_q, H(B\mid W), I(Y;B\mid W),
\dim \fib{w}, \dW, \Delta_{ij}(q))$ --- one record consolidating Rungs
5--9; the cached tension provably matches the bundle (\PROVEN).
\end{definition}

\begin{remark}[Why ``multi fractal workflow'']
The name is now redeemable in full: the pushforward of behavioral measure
onto the POWL hierarchy is (conjecturally) multifractal --- different
semantic freedoms (behavior, time, choice, linearization, slack, fluents)
scale with genuinely different exponents across the same hierarchy, and
the workflow's structure is read off from the \emph{spread} of
$\vec{D}_q$, not from any single dimension.
\end{remark}

% ================================================================
\section{Rung 10: Stratified differential geometry --- dimension loss}
\label{sec:rung10}

\subsection{Mathematics}
A \emph{manifold} looks locally like $\R^d$; its \emph{dimension} is $d$.
Spaces cut out by equalities $C(x) = 0$ and inequalities $G(x) \le 0$ with
discrete choices are not manifolds but \emph{stratified spaces}: unions of
manifolds (\emph{strata}) of varying dimension, glued along boundaries.
A \emph{submersion} has surjective differential; the preimage/fiber theorem
then makes fibers submanifolds of dimension
$\dim(\text{source}) - \dim(\text{target})$. Lipschitz maps cannot raise
dimension. The \emph{tangent space} at a point collects infinitesimal
directions; a \emph{null space} of a map collects directions it does not
see.

\subsection{In MFW}
\begin{definition}[Stratification of the phase space]\DEFINITION\;
Coordinates come in five families
$x = (\text{event occurrence}, \text{temporal}, \text{object-fluent},
\text{numeric}, \text{trajectory})$; discrete choice patterns
$\sigma \in \Sigma$ induce
$\PS = \bigcup_\sigma P_\sigma$ with covering law, each stratum carrying a
local dimension (continuous freedom: slack, numeric state).
\end{definition}

\begin{definition}[Dimension loss]\DEFINITION\;
$\Delta d_\tau(b) = d^{\mathrm{PDDL}}_{\mathrm{loc}}(b)
- d^{\mathrm{POWL}}_{\mathrm{loc}}(\tau(b)) \in \Z$. Strictly stronger
information than cardinality: an uncountable fiber (continuous slack) can
be low-dimensional; a finite fiber has dimension 0. Honest ledger: the
POWL-side local dimension is currently the placeholder $0$ (untimed
structural POWL), so the nonnegativity ``proof'' is vacuously easy; the
real statement needs a Lipschitz hypothesis (\CONJECTURAL).
``$\Delta d_\tau = 0 \iff$ local invertibility'' and ``fiber dimension
$=$ dimension loss under submersion'' are \BLOCKED: they require smooth
structure on $\PS$ and $\WS$ not yet formalized.
\end{definition}

\begin{definition}[Null perturbation space]\DEFINITION\;
$\mathrm{Null}_\tau(b) = \{\delta : \tau(b \oplus \delta) = \tau(b)\}$,
the lawful perturbations invisible to the compiler --- the fiber tangent
space at regular points. Proven basics: the trivial perturbation is null;
kind-restricted null spaces are subsets; null perturbations preserve
$\sim_\tau$ and stay in the fiber (\PROVEN). The kind decomposition
$\Delta d_\tau = \Delta d^{\mathrm{temp}} + \Delta d^{\mathrm{fluent}}
+ \Delta d^{\mathrm{num}} + \Delta d^{\mathrm{traj}}$ is \BLOCKED{}
(needs direct-sum tangent decomposition); per-kind dimensions are
placeholders. The measure link: if the source measure is $d$-dimensional
and fibers are $k$-dimensional,
$\mu(B_\varepsilon(w)) \sim \varepsilon^{d-k}
\mathrm{Vol}_k(\fib{w} \cap B_\varepsilon)$ --- dimension loss controls how
Rung 8's pushforward concentrates, feeding Rung 9's exponents.
\end{definition}

The physics-flavored open question of this rung, imported from
black-hole thermodynamics via the fiber-entropy scaling models
(Rung 5 $\times$ Rung 10): does $S_\tau(C)$ scale with component
\emph{volume} $|C|$ or \emph{boundary} $|\partial C|$
(\emph{boundary sufficiency hypothesis},
$\exists \beta > 0: |S_\tau(C) - \beta|\partial C|| \le \varepsilon$)?
An area law would mean the POWL interface variables suffice to reconstruct
interior behavioral information --- a holographic property of the compiler.
\DEFINITION{} (hypotheses stated; empirical status open).

% ================================================================
\section{Rung 11: Algebraic topology --- the shape of concurrency}
\label{sec:rung11}

\subsection{Mathematics}
An \emph{abstract simplicial complex} is a family of finite sets closed
under subsets --- a combinatorial space whose $k$-faces are $(k{+}1)$-element
sets. \emph{Simplicial homology} $H_k$ counts $k$-dimensional holes;
\emph{persistent homology} tracks holes across a filtration parameter.
A complex is \emph{contractible} when it has no holes at all.

\subsection{In MFW}
\begin{definition}[Concurrency complex]\DEFINITION\;
$K_\PT = \{A \subseteq E : A \text{ jointly concurrent}\}$. Downward
closure and $\emptyset \in K$ are \emph{admission hypotheses}, not free
(the audit removed sorry'd closure proofs --- they are false for arbitrary
predicates, e.g.\ ``$|A| = 2$''); with them, $K_\PT$ is a genuine
simplicial complex with dimension $\max_{A} |A| - 1$ (\PROVEN{}
construction given the hypotheses).
\end{definition}

The research program: $H_k(K_\PT)$ and its persistence as temporal or
constraint thresholds vary would measure the \emph{shape} of concurrency
(not just its width); the transformation
$\mathrm{PDDL} \to K_\PT \to \mathrm{POWL}$ may be selecting a hierarchical
representation of this complex. \emph{Falsifier recorded in-source}: the
thesis dies if $K_\PT$ is trivially contractible for all admitted planning
theories. All homological claims: \CONJECTURAL.

% ================================================================
\section{Rung 12: The generator picture --- presentation of the kernel}
\label{sec:rung12}

\subsection{Mathematics}
An equivalence relation can be \emph{presented} by generators: primitive
moves whose reflexive--symmetric--transitive closure it is (as groups are
presented by generators and relations). If two generator types are
independent (neither's moves are expressible by the other's), quantities
counting them are independent coordinates.

\subsection{In MFW}
\begin{definition}[Kernel generators and paths]\DEFINITION\;
$g \in \{\mathrm{traceSwap}, \mathrm{choiceNormalization},
\mathrm{hierarchyCollapse}, \mathrm{temporalAbstraction}\}$; a
\emph{kernel path} $b_1 \to^{g_1} \cdots \to^{g_n} b_2$ carries
generators, intermediates, and endpoint proofs; $\kappa(\mathrm{path})$
counts each type.
\end{definition}

\begin{conjecture}[Kernel generated]\CONJECTURAL\;
$b_1 \kerEq b_2 \iff \exists$ kernel path. If it holds, the generator
types determine which erasure modes are independent, hence \emph{which
spectra exist}: one $D_q$ per independent generator --- the proof-first
derivation of Rung 9's basis, replacing enumeration of measure kinds by
presentation of the kernel. This is the deepest expression of the paper's
organizing principle: \emph{first characterize what the map forgets; all
structure follows.}
\end{conjecture}

% ================================================================
\section{Rung 13: Governance mathematics --- the control plane}
\label{sec:rung13}

The remaining formal modules are interface-level specifications
(\texttt{opaque} carriers, \texttt{Prop}-valued laws) that wrap the
analytic core in an engineering discipline; their mathematics is modest
but their \emph{shape} is load-bearing:

\begin{itemize}[nosep]
  \item \textbf{Manufacture / BRCE} (\S\S1--6, 15): manufacturing law
  $\mathcal{L} \subseteq \mathcal{A} \times \mathcal{O}$, receipts, and
  the invariant $\forall a: \mathrm{isReceipted}(a)$ (zero unreceipted
  actuation); Rice containment scopes verification to an admitted semantic
  domain --- a nod to Rice's theorem that arbitrary semantic properties are
  undecidable. \DEFINITION
  \item \textbf{Falsification} (\S38): status
  $\{\mathrm{Admitted}, \mathrm{Revoked}\}$;
  $\mathrm{RevCond} = (\mathrm{Admitted} \wedge F(T,E))$;
  $\mathrm{FalsLoop} = (\mathrm{RevCond} \implies
  \mathrm{status}(\mathrm{revoke}) = \mathrm{Revoked})$ --- Popper as an
  interface. Composition proven from hypotheses in tests. \DEFINITION
  \item \textbf{Explore/Exploit} (\S\S7--12, 35): exploitation factors
  through contracts (a congruence-shaped separation law). \DEFINITION
  \item \textbf{Search portfolio} (\S\S31--35): an exact completeness rail
  $F_0$ plus q-lens-parameterized exploitation rails; conjecture:
  completeness $+$ persistent service $+$ noninterference $\implies$
  portfolio completeness. \CONJECTURAL
  \item \textbf{Compiler pipeline} (\S\S36--37): the cascade
  $\mathrm{Coords} \to \mathrm{TTL} \to \mathrm{Tera} \to
  \mathrm{RustExe}$ with stage equalities carried as proof fields and the
  bound $\mathrm{multDepth} \le \mathrm{time}$. \DEFINITION
  \item \textbf{Ledger} (\S\S39--40): claims as a DAG with derivations in
  topological order --- \emph{known stub}:
  $\mathrm{validTopologicalSort} := \mathrm{True}$; the check validates
  nothing yet. \DEFINITION{} (placeholder)
\end{itemize}

% ================================================================
\section{The whole ladder at a glance}
\label{sec:summary}

\begin{center}
\small
\begin{tabular}{@{}llll@{}}
\toprule
Rung & Mathematics & MFW object & Peak standing \\
\midrule
0 & sets, partitions, counting & fibers; $p(w)$ & \PROVEN{} (disjointness) \\
1 & equivalence, quotients, posets & $\kerEq$, $\sim_I$, $[b]_I$, crown & \CONJECTURAL{} (crown) \\
2 & graphs, Petri/WF-nets & POWL v2; separable conversion & \PROVEN{} (Thms 5.5/5.6) \\
3 & linear extensions, Dilworth & $\Hser$, $R^*$, $H_T$ gap, $C_E$ & \DEFINITION \\
4 & dependent type theory & proof-carrying behaviors; \texttt{rfl} & \PROVEN{} (non-circularity) \\
5 & Shannon entropy, sufficiency & $S_\tau$, $H(B|W)$, $\eta$, collapse & \CONJECTURAL{} (collapse) \\
6 & sigma-algebras, vector spaces & $\mathcal{O}_\tau$, $F_\tau$, basis & \PROVEN{} (algebra of $F_\tau$) \\
7 & escort distributions & q-lens, ratio law & \PROVEN \\
8 & measures, optimal transport & $\tau_*\nu$, $\nu_w$, $\dW = W_p$ & \DEFINITION{} (opaque) \\
9 & R\'enyi dimensions, multifractals & $\vec D_q$, $\Delta_{ij}(q)$, $Z(q,k)$ & \PROVEN{} (tension algebra) \\
10 & stratified geometry & $\Delta d_\tau$, $\mathrm{Null}_\tau$, area law & \BLOCKED{} (smoothness) \\
11 & simplicial complexes, homology & $K_\PT$, persistence & \CONJECTURAL \\
12 & presentations by generators & kernel paths, spectrum basis & \CONJECTURAL \\
13 & specification logic & BRCE, falsification, ledger & \DEFINITION \\
\bottomrule
\end{tabular}
\end{center}

Reading the table column-wise gives the paper's three theses. First
(column 2): no rung is decorative --- each field enters because a specific
question about $\tau$ cannot be asked without it. Second (column 3): every
object at every level is derived from, or aimed at, the kernel; the crown
conjecture is the single point of failure and the single point of
unification. Third (column 4): the formalization keeps an exact ledger of
what is settled. The proven layer (partitions, trace algebra, conversion
theorems, non-circularity, observable algebra, q-lens, tension algebra) is
already enough to build and audit a compiler; the conjectural layer (crown,
collapse, basis, generation) says what such a compiler would \emph{mean};
the blocked layer (smooth structure) marks where new formalized mathematics
--- not new ideas --- is the bottleneck.

\begin{thebibliography}{9}
\bibitem{mfw} \texttt{ProcInt.MFW}, Lean~4 formalization, toolchain
v4.31.0, Mathlib pinned \texttt{fabf563a}; companion papers
\texttt{mfw-notation.tex} (full notation) and \texttt{mfw-tutorial.tex}
(the conversion question).
\bibitem{kourani2026} H.~Kourani, G.~Park, W.~M.~P.~van der Aalst,
\emph{Hierarchical Decomposition of Separable Workflow-Nets},
arXiv:2602.15739v3, 2026.
\bibitem{powl} H.~Kourani, S.~J.~van Zelst, \emph{POWL}, BPM 2023.
\bibitem{traces} A.~Mazurkiewicz 1977; V.~Diekert, G.~Rozenberg,
\emph{The Book of Traces}, 1995.
\bibitem{villani} C.~Villani, \emph{Optimal Transport: Old and New},
Springer, 2009.
\bibitem{efp} P.~Komiske, E.~Metodiev, J.~Thaler, \emph{Energy Flow
Polynomials}, JHEP 2018.
\bibitem{pddl} M.~Fox, D.~Long, PDDL 2.1/3.x.
\end{thebibliography}

\end{document}
